\documentclass[12pt]{article}
\usepackage[utf8]{inputenc}
\usepackage[spanish]{babel}
\usepackage{amsmath}
\usepackage{amsfonts}
\usepackage{amssymb}
\usepackage{geometry}
\geometry{letterpaper, margin=2.5cm}

\begin{document}

\begin{center}
    \Large \textbf{Evaluación de Examen 1: Unidad I} \\
    \large Física para Ciencias de la Salud (005-1713) \\
    \vspace{0.5cm}
    \textbf{Estudiante:} Luisannys García \quad \textbf{C.I.:} 33.141.963
    \vspace{0.3cm} \\
    \large \textbf{Calificación Final: 9/10 puntos}
\end{center}

\vspace{0.5cm}

\section*{Análisis de Resultados}

\subsection*{1. Ángulo plano (Conversión a radianes)}
\textbf{Procedimiento y resultado:} Plantea de forma sumamente ordenada una regla de tres basándose en la equivalencia $180^\circ \rightarrow \pi$ rad. Desarrolla la multiplicación y simplifica correctamente la fracción, obteniendo el valor exacto de $\frac{5\pi}{12}$ rad. \\
\textbf{Veredicto:} \textbf{Correcto} (2/2 pts).

\subsection*{2. Sistemas de coordenadas (Longitud del hueso)}
\textbf{Procedimiento y resultado:} Excelente resolución. La estudiante declara explícitamente las coordenadas y la fórmula de distancia entre dos puntos antes de sustituir. Realiza los cálculos de diferencias y potencias paso a paso ($\sqrt{36+64} = \sqrt{100}$) y obtiene impecablemente la longitud de $10$ cm. \\
\textbf{Veredicto:} \textbf{Correcto} (2/2 pts).

\subsection*{3. Vectores (Fuerza resultante)}
\textbf{Procedimiento y resultado:} Muestra un gran orden al desglosar el cálculo en sus componentes rectangulares. Suma los valores en el eje x ($45 - 15 = 30$) y en el eje y ($25 + 35 = 60$) de forma aislada para evitar errores. Escribe el vector de fuerza resultante de manera formal y con las unidades correctas: $(30\hat{i} + 60\hat{j})$ N. \\
\textbf{Veredicto:} \textbf{Correcto} (2/2 pts).

\subsection*{4. Potencias de diez (Notación científica y prefijo)}
\textbf{Procedimiento y resultado:} Transforma el número decimal a notación científica estándar de forma matemáticamente correcta ($1,25 \times 10^{-5}$ m). Al igual que muchos de sus compañeros, omitió nombrar el prefijo del S.I. (micrómetros) que solicitaba la segunda parte del enunciado. \\
\textbf{Veredicto:} \textbf{Incompleto / Parcialmente correcto} (1/2 pts).

\subsection*{5. Sistemas de unidades (Conversión de densidad)}
\textbf{Procedimiento y resultado:} Aunque no desarrolla la cadena de factores de conversión factor a factor, demuestra conocer de antemano el factor de conversión global, estableciendo que $1 \text{ g/cm}^3 = 1000 \text{ kg/m}^3$. Ejecuta la multiplicación correctamente ($1,03 \times 1000$) llegando a un resultado numérico y unidades exactas: $1030 \text{ Kg/m}^3$. \\
\textbf{Veredicto:} \textbf{Correcto} (2/2 pts).

\vspace{0.5cm}
\section*{Comentarios Generales y Sugerencias}
¡Excelente examen! La estudiante destaca por su impecable orden, uso del espacio y claridad al estructurar sus procedimientos matemáticos (fórmulas, sustitución y resultados resaltados).
\begin{itemize}
    \item Se recomienda prestar un poco más de atención a la lectura detallada de los problemas para evitar omitir partes de la respuesta, como ocurrió al no escribir el nombre del prefijo en la Pregunta 4.
\end{itemize}

\end{document}